\documentclass[12pt,a4paper]{article}
\usepackage[utf8]{inputenc}
\usepackage[spanish]{babel}
\usepackage{amsmath}
\usepackage{amssymb}
\usepackage{listings}
\usepackage{enumitem}
\usepackage{geometry}
\geometry{margin=2.5cm}
\usepackage[utf8]{inputenc} % allow utf-8 input

\usepackage{listings}
\usepackage{xcolor}

% Define colors
\definecolor{codegreen}{rgb}{0,0.6,0}
\definecolor{codegray}{rgb}{0.5,0.5,0.5}
\definecolor{codepurple}{rgb}{0.58,0,0.82}
\definecolor{backcolour}{rgb}{0.95,0.95,0.92}

% Configure listings
\lstdefinestyle{mystyle}{
  backgroundcolor=\color{backcolour}, 
  commentstyle=\color{codegreen},
  keywordstyle=\color{magenta},
  numberstyle=\tiny\color{codegray},
  stringstyle=\color{codepurple},
  basicstyle=\ttfamily\footnotesize,
  breakatwhitespace=false,   
  breaklines=true,       
  captionpos=b,        
  keepspaces=true,       
  numbers=left,        
  numbersep=5pt,      
  showspaces=false,      
  showstringspaces=false,
  showtabs=false,      
  tabsize=2
}

\lstset{style=mystyle}


\setlength{\parindent}{0pt} % Eliminar sangría de primera línea en todos los párrafos
% add space between paragraphs
\setlength{\parskip}{1em}

% agregar carpeta imagenes
\usepackage{graphicx}
\graphicspath{ {../imagenes/} }

\title{Aprendizaje Reforzado: Guía 4: Métodos Monte Carlo}
\author{José Saint Germain}

\begin{document}

\maketitle

\section{Ejercicio 1}

Considere un proceso de decisión de Markov (MDP) con dos estados: uno terminal y otro no-terminal.
Existe sólo una acción que lleva del estado no-terminal al estado terminal con probabilidad 1-p y 
del estado no-terminal a sí mismo con probabilidad p. La recompensa es +1 en todas las
transiciones y el factor de descuento es $\gamma = 1$. Suponga que observa un episodio con
10 iteraciones y un retorno de 10.

\begin{enumerate}
    \item ¿Cuáles son las estimaciones Monte Carlo de primer-visita y de cada-visita
    del valor del estado no-terminal basadas en ese episodio?
    \item Compare los valores de estado obtenidos con el teórico p=0.9. Saque conclusiones.
\end{enumerate}

\subsection{Resolución}

\subsubsection{Parte 1}

En el método de primera visita, solo consideramos la primera vez que visitamos el estado no-terminal.

\begin{itemize}
    \item \textbf{Primera visita al estado no-terminal:} tiempo $t=0$
    \item \textbf{Retorno desde $t=0$:} Como $\gamma = 1$, el retorno es la suma de todas las recompensas futuras:
    
    $$G_0 = r_1 + r_2 + r_3 + \cdots + r_{10} = 1 + 1 + 1 + \cdots + 1 = 10$$
\end{itemize}

En el método de cada visita, consideramos todas las veces que visitamos el estado no-terminal.

\begin{itemize}
    \item \textbf{Visita en $t=0$:} $G_0 = 10$
    \item \textbf{Visita en $t=1$:} $G_1 = 9$
    \item \textbf{Visita en $t=2$:} $G_2 = 8$
    \item \textbf{Visita en $t=3$:} $G_3 = 7$
    \item \textbf{Visita en $t=4$:} $G_4 = 6$
    \item \textbf{Visita en $t=5$:} $G_5 = 5$
    \item \textbf{Visita en $t=6$:} $G_6 = 4$
    \item \textbf{Visita en $t=7$:} $G_7 = 3$
    \item \textbf{Visita en $t=8$:} $G_8 = 2$
    \item \textbf{Visita en $t=9$:} $G_9 = 1$
\end{itemize}

\textbf{Estimación de cada visita:}
$$V_{\text{every}}(S_0) = \frac{1}{10}\sum_{i=0}^{9} G_i = \frac{10 + 9 + 8 + 7 + 6 + 5 + 4 + 3 + 2 + 1}{10} = \frac{55}{10} = 5.5$$

\subsubsection{Parte 2}

Para calcular el valor teórico del estado no-terminal con $p = 0.9$, usamos la ecuación de Bellman:

$$V(S_0) = \mathbb{E}[R + \gamma V(S')]$$

donde:
\begin{itemize}
    \item Con probabilidad $1-p = 0.1$: transición al terminal, $R=1$, $V(S_T) = 0$
    \item Con probabilidad $p = 0.9$: transición a sí mismo, $R=1$, $V(S_0)$
\end{itemize}

La ecuación queda:
$$V(S_0) = (1-p) \cdot (1 + \gamma \cdot 0) + p \cdot (1 + \gamma \cdot V(S_0))$$
$$= (1-p) \cdot 1 + p \cdot (1 + 1 \cdot V(S_0))$$
$$= (1-p) + p + pV(S_0)$$

Resolviendo para $V(S_0)$:
$$V(S_0) = 1 + pV(S_0)$$
$$V(S_0) - pV(S_0) = 1$$
$$(1-p)V(S_0) = 1$$
$$V(S_0) = \frac{1}{1-p}$$

Para $p = 0.9$:
$$V_{\text{teórico}}(S_0) = \frac{1}{1-0.9} = \frac{1}{0.1} = 10$$


En este episodio particular, la estimación de \textit{primera-visita} coincide 
exactamente con el valor verdadero (10). Esto tiene sentido: con $\gamma = 1$ y 
recompensa $+1$ en cada transición, el retorno desde la primera visita es simplemente 
el número total de pasos del episodio, que fue 10. La estimación de \textit{cada-visita} 
es menor (5.5) porque promedia retornos decrecientes a lo largo del episodio.

\section{Ejercicio 2}

Demostrar que la política $\epsilon-Greedy$, definida de la siguiente manera

\begin{equation*}
    \pi(a|s) = 
    \begin{cases} 
      1 - \epsilon + \frac{\epsilon}{|A(s)|} & \text{si } a = a^* \\
      \frac{\epsilon}{|A(s)|} & \text{si } a \neq a^*
    \end{cases}
\end{equation*}

es una distribución de probabilidad válida, donde $|A(s)|$ es el número de acciones
posibles para el estado $s,\epsilon<1$ es un número positivo pequeño, y $a^*$ es la
acción óptima (decisión greedy) para el estado s ¿Hay que pedir alguna condición sobre
$\epsilon$?

\subsection{Resolución}

Para que $\pi(a|s)$ sea una distribución de probabilidad válida, debe cumplir dos condiciones:
\begin{enumerate}
    \item Normalización: La suma de las probabilidades de todas las acciones posibles debe ser igual a 1, es decir, $\sum_{a \in A(s)} \pi(a|s) = 1$.
    \item No negatividad: $\pi(a|s) \geq 0$ para todas las acciones $a$ en el conjunto de acciones $A(s)$.
\end{enumerate}

\subsubsection{No negatividad}

Para $A\neq a^* \pi(a|s) = \frac{\epsilon}{|A(s)|} \geq 0$.
Esto se puede demostrar que es mayor o igual a 0 si $\epsilon \geq 0$.

Si $A=a^* \pi(a|s) = 1-\epsilon + \frac{\epsilon}{|A(s)|} \geq 0$

Se puede reescribir como $1-\epsilon(1-\frac{1}{|A(s)|})$

$1-\epsilon(1-\frac{1}{|A(s)|}) \geq 0$, entonces la expresión va a ser mayor o
igual a 0 si $1\geq \epsilon(1-\frac{1}{|A(s)|})$.

De aquí surge una condición necesaria para $\epsilon$: $\epsilon \leq \frac{1}{|A(s)|-1}$.

\subsubsection{Normalización}

$$\sum_a \pi(a|s) = \pi(a^*|s) + \sum_{a\neq a^*} \pi(a|s)$$

Remplazando los valores de $\pi(a|s)$:

$$= 1 - \epsilon + \frac{\epsilon}{|A(s)|} + \sum_{a\neq a^*} \frac{\epsilon}{|A(s)|}$$

Remplazamos por el número de acciones menos 1 (la acción óptima):

$$= 1 - \epsilon + \frac{\epsilon}{|A(s)|} + (|A(s)|-1) \frac{\epsilon}{|A(s)|}$$

$$= 1 - \epsilon + \frac{\epsilon}{|A(s)|} + \epsilon - \frac{\epsilon}{|A(s)|}$$

$$= 1 - \epsilon + \epsilon = 1$$

Por lo tanto, se cumple la condición de normalización.

\section{Ejercicio 3}

Se dice que una política $\pi$ es $\epsilon-soft$ si 
$\pi(a|s) \geq \frac{\epsilon}{|A(s)|} \forall a \neq a^* y \forall s$.
Demostrar que la política $\epsilon-Greedy \pi'(a|s)\geq\frac{\epsilon}{|A(s)|}$
(definida en el ejercicio 1) es igual o mejor que cualquier política 
$\epsilon-soft$, es decir, $v_{\pi'}(s) \geq v_{\pi}(s) \forall s$.

\subsection{Resolución}

Con la información recibida, tenemos que demostrar (por el teorema de la 
mejora de la política) que:
$$q_\pi(s,\pi'(s)) \geq \nu_\pi(s)$$

Sabemos que:
$$q_{\pi}(s,\pi'(s)) = E[G_t|S_t=s,A_t=\pi'(s)]$$
como $A_t=\pi'(s)$ es determinístico, lo podemos reescribir:
$$ q_{\pi}(s,\pi'(s)) = E_\pi[G_t|S_t=s] = \sum_g gp(g|s)$$
hacemos la probabilidad marginal:
$$= \sum_g \sum_a p(g,a|s)$$
aplicamos el teorema de Bayes:
$$= \sum_g \sum_a p(g|s,a)p(a|s)$$
sabiendo que $p(a|s) = \pi'(a|s)$ y $\sum_g p(g|s,a) = q(s,a)$:
$$q_{\pi}(s,\pi'(s))= \sum_a \pi'(a|s)q(s,a)$$

esta sumatoria se puede separar en la acción óptima y las demás:
$$= \pi'(a^*|s)q(s,a^*) + \sum_{a\neq a^*} \pi'(a|s)q(s,a)$$
remplazamos los valores de $\pi'(a|s)$:
$$= (1-\epsilon + \frac{\epsilon}{|A(s)|})q_\pi(s,a^*) + \sum_{a\neq a^*} \frac{\epsilon}{|A(s)|}q_\pi(s,a)$$
$$= (1-\epsilon)q_\pi(s,a^*) + \frac{\epsilon}{|A(s)|} \sum_{a} q_\pi(s,a)$$
distribuyendo $q_\pi(s,a^*)$:
$$ q_\pi(s,a^*) (1-\epsilon) + q_\pi(s,a^*) \frac{\epsilon}{|A(s)|} + \frac{\epsilon}{|A(s)|} \sum_{a\neq a^*} q_\pi(s,a)$$
junto las sumatorias $\forall a$:
$$= q_\pi(s,a^*) (1-\epsilon) + \frac{\epsilon}{|A(s)|} \sum_{a} q_\pi(s,a)$$
remplazo $q_\pi(s,a)$ por:
$$= q_\pi(s,a^*) = \sum_a c(a) q_\pi(s,a) \text{ si} \sum_a c(a) = 1 $$

Asumimos que $q \sum_a c(a) q_\pi(s,a) \leq q_\pi(s,a^*)$

$$ q_\pi(s,a^*) \geq \frac{\epsilon}{|A(s)|}\sum_{|A(s)|}q_\pi(s,a)+(1-\epsilon)\sum_a c(a) q_\pi(s,a)$$

Elijo unos coeficientes definidos como:
$$ c(a) = \frac{\pi(a|s) - \frac{\epsilon}{|A(s)}}{1-\epsilon}$$
sabiendo que la sumatoria de todos los coeficientes es 1:
$$ \sum_a c(a) = \sum_a \frac{\pi(a|s) - \frac{\epsilon}{|A(s)}}{1-\epsilon} = \frac{1 - \epsilon}{1-\epsilon} = 1$$
remplazando c(a) en la desigualdad:

$$ q_\pi(s,a^*) \geq \frac{\epsilon}{|A(s)|}\sum_{|A(s)|}q_\pi(s,a)+\sum_a \pi(a|s)q_\pi(s,a) - \frac{\epsilon}{|A(s)|} \sum_a q_\pi(s,a)$$
cancelo los términos semejantes:
$$ q_\pi(s,a^*) \geq \sum_a \pi(a|s)q_\pi(s,a) \stackrel{\text{definición}}{=} v_\pi(s)$$
entonces,
$$ q_\pi(s,\pi'(s)) \geq v_\pi(s)$$
que es equivalente (por teorema de la mejora de la política) a:
$$ \nu_{\pi'}(s) \geq \nu_\pi(s)$$

\section{Ejercicio 4}

Dada una trayectoria de acciones y estados $A_t,S_{t+1},A_{t+1},\dots,S_T$ en un
proceso de decisión markoviano (MDP) con política $\pi(a|s)$, demostrar que la 
probabilidad conjunta de esa trayectoria se puede escribir como:

$$P[S_t,A_t,S_{t+1},A_{t+1},\dots,S_T] = \prod_{k=t}^{T-1} \pi(A_k|S_k)p(S_{k+1}|S_k,A_k)$$

Nota: considere $P[S_t]=1$.

\subsection{Resolución}

Por la regla de la cadena:
\begin{align*}
P[S_t, A_t, S_{t+1}, A_{t+1}, \ldots, S_T] = P[S_t] &\cdot P[A_t|S_t] \\
&\cdot P[S_{t+1}|S_t, A_t] \\
&\cdot P[A_{t+1}|S_t, A_t, S_{t+1}] \\
&\cdot P[S_{t+2}|S_t, A_t, S_{t+1}, A_{t+1}] \\
&\cdots
\end{align*}

En un MDP el siguiente estado depende solo del estado y acción actuales:
$$P[S_{k+1}|S_t, A_t, \ldots, S_k, A_k] = P(S_{k+1}|S_k, A_k)$$
    
Y por política estocástica, la acción depende solo del estado actual:
$$P[A_k|S_t, A_t, \ldots, S_k] = \pi(A_k|S_k)$$

Aplicando estas propiedades:
$$P[S_t, A_t, S_{t+1}, \ldots, S_T] = P[S_t] \prod_{k=t}^{T-1} \pi(A_k|S_k) \cdot P(S_{k+1}|S_k, A_k)$$

Dado que $P[S_t] = 1$, obtenemos:

$$P[S_t, A_t, S_{t+1}, A_{t+1}, \ldots, S_T] = \prod_{k=t}^{T-1} \pi(A_k|S_k)P(S_{k+1}|S_k, A_k)$$


\section{Ejercicio 5}

Usando el resultado del ejercicio 4, demostrar que el importance-sampling ratio 
correspondiente a la aplicación del método off-policy es:

$$ \rho_{t:T-1} = \prod_{k=t}^{T-1} \frac{\pi(A_k|S_k)}{b(A_k|S_k)} $$

\subsection{Resolución}

Sabiendo que:
$$P_{\pi}[S_t, A_t, S_{t+1}, A_{t+1}, \ldots, S_T] = \prod_{k=t}^{T-1} \pi(A_k|S_k)P(S_{k+1}|S_k, A_k)
$$

De manera análoga, la probabilidad de la misma trayectoria bajo la política de comportamiento $b$ es:
$$P_{b}[S_t, A_t, S_{t+1}, A_{t+1}, \ldots, S_T] = \prod_{k=t}^{T-1} b(A_k|S_k)P(S_{k+1}|S_k, A_k)$$

El \textit{importance-sampling ratio} se define como el cociente entre la probabilidad de la trayectoria bajo la política objetivo y la probabilidad bajo la política de comportamiento:

$$\rho_{t:T-1} = \frac{P_{\pi}[S_t, A_t, S_{t+1}, A_{t+1}, \ldots, S_T]}{P_{b}[S_t, A_t, S_{t+1}, A_{t+1}, \ldots, S_T]} \\[10pt]$$
$$= \frac{\prod_{k=t}^{T-1} \pi(A_k|S_k)P(S_{k+1}|S_k, A_k)}{\prod_{k=t}^{T-1} b(A_k|S_k)P(S_{k+1}|S_k, A_k)} \\[10pt]$$
$$= \frac{\prod_{k=t}^{T-1} \pi(A_k|S_k) \cdot \prod_{k=t}^{T-1} P(S_{k+1}|S_k, A_k)}{\prod_{k=t}^{T-1} b(A_k|S_k) \cdot \prod_{k=t}^{T-1} P(S_{k+1}|S_k, A_k)}$$

simplificamos la fracción:

$$= \prod_{k=t}^{T-1} \frac{\pi(A_k|S_k)}{b(A_k|S_k)}$$

\section{Ejercicio 6}

Demostrar que la fórmula de la implementación incremental de un promedio ponderado, es decir,
el promedio ponderado

$$ V_n = \frac{\sum_{k=1}^{n-1} W_k G_k}{\sum_{k=1}^{n-1} W_k} $$

puede calcularse incrementalmente con la fórmula:

$$ V_{n+1} = V_n + \frac{W_n}{C_n}[G_n - V_n] $$

con $C_{n+1} = C_n + W_{n+1}$.

\subsection{Resolución}

Sea $C_n = \sum_{k=1}^{n} W_k$ el acumulador de pesos hasta el paso $n$. Entonces:
$$V_n = \frac{\sum_{k=1}^{n} W_k G_k}{C_n}$$

Por definición:
$$V_{n+1} = \frac{\sum_{k=1}^{n+1} W_k G_k}{C_{n+1}} = \frac{\sum_{k=1}^{n+1} W_k G_k}{C_n + W_{n+1}}$$

Podemos escribir:
$$V_{n+1} = \frac{\sum_{k=1}^{n} W_k G_k + W_{n+1}G_{n+1}}{C_n + W_{n+1}}$$

Sabemos que $\sum_{k=1}^{n} W_k G_k = V_n \cdot C_n$, entonces:
$$V_{n+1} = \frac{V_n C_n + W_{n+1}G_{n+1}}{C_n + W_{n+1}}$$

Como $C_{n+1} = C_n + W_{n+1}$:
$$V_{n+1} = \frac{V_n C_n + W_{n+1}G_{n+1}}{C_{n+1}} \\[8pt]$$
$$= \frac{V_n C_n}{C_{n+1}} + \frac{W_{n+1}G_{n+1}}{C_{n+1}}$$

Podemos reescribir el primer término:
$$V_{n+1} = \frac{V_n C_n}{C_{n+1}} + \frac{W_{n+1}G_{n+1}}{C_{n+1}} \\[8pt]$$
$$= \frac{V_n(C_{n+1} - W_{n+1})}{C_{n+1}} + \frac{W_{n+1}G_{n+1}}{C_{n+1}} \\[8pt]$$
$$= \frac{V_n C_{n+1} - V_n W_{n+1}}{C_{n+1}} + \frac{W_{n+1}G_{n+1}}{C_{n+1}} \\[8pt]$$
$$= V_n - \frac{V_n W_{n+1}}{C_{n+1}} + \frac{W_{n+1}G_{n+1}}{C_{n+1}}$$

$$V_{n+1} = V_n + \frac{W_{n+1}G_{n+1} - V_n W_{n+1}}{C_{n+1}} \\$$
$$= V_n + \frac{W_{n+1}}{C_{n+1}}(G_{n+1} - V_n)$$

Si reemplazamos $n+1$ por $n$ en la fórmula final (es decir, consideramos la actualización del paso $n-1$ al paso $n$), obtenemos:

$$V_{n+1} = V_n + \frac{W_n}{C_n}[G_n - V_n]$$

\end{document}